\section{Proofs}\label{sec:suppMaterial}

\subsection{Proof of \lemref{lem:analyticalSolutionK=2}}\label{sec:proof_lem:analyticalSolutionK=2}
\begin{proof}
Our starting point is the following lemma, which makes further analysis easier.
\begin{lemma}\label{lem:replaceSamplesWithMean}
Assume an optimal solution $\vmu^*$ of \eqref{eq:segmentationObjective_mu} is given, and therefore we also know to which
segment each data sample belongs. If we replace all samples in a segment with the mean of these samples, the optimal
solution $\vmu^*$ will not change.
\end{lemma}
The proof appears in \secref{sec:proof_lem:replaceSamplesWithMean}. % of the appendix. % in the supplementary material.
We now analyze the transition of the solution to \eqref{eq:segmentationObjective_mu} from $K=1$ to $K=2$ segments.
During the analysis we use the fact that the solution path is continuous in $\lambda$, as was shown previously in
another context by \cite{Chi:2013arXiv}. We denote by $\lambda^*$ the value of $\lambda$ at the splitting point, and we
assume that for the two segments solution we have $n_1$ samples in the first segment and $n_2=n-n_1$ samples in the second
segment. We denote the means of the two segments by $\bvxi{1}$ and $\bvxi{2}$. \lemref{lem:replaceSamplesWithMean} allows
us to replace samples in a segment with the mean of the segment, without changing the optimal solution $\vmu^*$.
This means that for $K=2$ all analysis is taking place on the line connecting $\bvxi{1}$ and $\bvxi{2}$ and is
therefore essentially one dimensional. For $K=1$ the regularization term vanishes, so the solution which we denote
by $\vmu$ is simply the mean of the whole data set, $\vmu=\bvxi{1}+\alpha_0(\bvxi{2}-\bvxi{1})$,
where $\alpha_0=n_2/n$. For $K=2$, we denote the solution by $\vmui{1}$ and $\vmui{2}$,
and we parameterize $\vmui{1}$ by $\vmui{1}=\bvxi{1}+\alpha\paren{\bvxi{2}-\bvxi{1}}$, for some $\alpha$.
We note that $\alpha\leq\alpha_0$, since we know that $\vmui{1}$ is closer to $\bvxi{1}$ than $\vmu$ is.
The parametrization for $\vmui{2}$ is therefore $\vmui{2}=\vxi{1}+\paren{1-\frac{n_1}{n_2}\alpha}\paren{\bvxi{2}-\bvxi{1}}$
\footnote{This can be derived either directly by requiring that the derivative of \eqref{eq:segmentationObjective_mu} for $K=2$
equals zeros, or by noting that the `center of mass' of $\vmui{1}$ and $\vmui{2}$ is just $\vmu$. Both approaches gives the
same equation, namely $n_1\vmui{1}+n_2\vmui{2}=n\vmu$.}.
In order to find $\alpha$, we look for a minimum of the objective $h$ for $K=2$:
\begin{align}\label{eq:objK=2}
    h =& \frac{n_1}{2}\normt{\bvxi{1}-\vmui{1}}+\frac{n_2}{2}\normt{\bvxi{2}-\vmui{2}}\\
    &+\lambda w_{n_1}\vnorm{\vmui{2}-\vmui{1}}_2.\nonumber
\end{align}
Plugging $\vmui{1}$ and $\vmui{2}$ as parameterized by $\alpha$ into \eqref{eq:objK=2} and looking for the minimum, we get
\begin{align*}%\label{eq:objOfAlpha}
\alpha & = \argmin{\alpha'}\braces{y^2\alpha'^2{nn_1}/{2n_2}-\lambda w_{n_1} y\alpha'{n}{n_2}+\lambda w_{n_1} y}\\
& = {\lambda w_{n_1}}/{n_1y},\nonumber
\end{align*}
where we define $y\triangleq \vnorm{\bvxi{2}-\bvxi{1}}$. In order to find $\lambda^*$, we require that the objective for $K=1$
and $K=2$ has the same value at the splitting point, where $\lambda=\lambda^*$. This requirement is equivalent to the requirement
that $\alpha=\alpha_0$, since at the splitting point we have $\vmu=\vmui{1}=\vmui{2}$. This leads to the solution for $\lambda^*$,
where we explicitly include the dependence on $n_1$, to emphasize that this is the solution provided that $n_1$ is known:
\begin{equation}\label{eq:lambdaStar_n1}
    \lambda^*(n_1)=\frac{1}{w_{n_1}}\frac{n_1(n-n_1)}{n}\vnorm{\bvxi{2}(n_1)-\bvxi{1}(n_1)}_2.
\end{equation}

%--- IK: once we moved this to the appendix, the following paragraph is not really needed for the proof
%Using these results, it is straightforward to calculate from \eqref{eq:objK=2} the derivative of the objective with respect
%to $\lambda$, which is found to be $\frac{dh}{d\lambda}=-\lambda\frac{n}{n_1n_2}+y$. Using \eqref{eq:lambdaStar_n1}
%we see that for $\lambda\to\lambda^*$ from below, the derivative is zero. Since for $\lambda>\lambda^*$ the derivative is zero
%as well, as the regularization term vanishes and the solution does not depend on $\lambda$, we conclude that the first
%derivative $\frac{dh}{d\lambda}$ is continuous at $\lambda^*$, but not smooth. Such a transition is often referred to as a
%second-order phase transition.

In order to find the actual splitting point $n^*$, we note that the split into two segments occurs as $\lambda$ is decreased
from $\lambda>\lambda^*$ to $\lambda<\lambda^*$, so maximizing $\lambda^*(n_1)$ over $n_1$ gives the splitting point $n^*$:
%\begin{align}\label{eq:lambdaStar}
%    \lambda^*&=\max_{n_1}\lambda^*(n_1)\nonumber\\
%    &=\max_{n_1}\frac{n_1\paren{n-n_1}}{n}\vnorm{\bvxi{2}\paren{n_1}-\bvxi{1}\paren{n_1}}_2.
%\end{align}
\begin{align*}%\label{eq:lambdaStar}
    n^*\paren{\vx} &= \argmax{n_1}g\paren{n_1,\vx},\\ %     \textrm{ and }
    \lambda^*\paren{\vx} &= g\paren{n^*,\vx},\nonumber \\
    \textrm{ where: }\\
    g\paren{n_1,\vx} &= \bbraces{\frac{1}{w_{n_1}}\frac{n_1(n-n_1)}{n}\vnorm{\bvxi{2}(n_1)-\bvxi{1}(n_1)}_2}.\nonumber
\end{align*}
\end{proof}\QED
%\newpage
\subsection{Proof of \lemref{lem:weighted_solution_l0_l1}}\label{sec:proof_lem:weighted_solution_l0_l1}
\begin{proof}
For $K=2$ the unrelaxed optimization problem \eqref{eq:segmentationObjective_unrelaxed_K=2} is equivalent to
%is given by
%\begin{align*}
%    \argmin{\vmu}&\braces{\sum_{i=1}^n\normt{\vxii-\vmuii}},\\
%    \mathrm{s.t.}\ &\sum_i\mb{1}\braces{\vnorm{\vmui{i+1}-\vmuii}_2>0}=1~,
%\end{align*}
%which is equivalent to minimizing
\begin{align*}
    \argmin{n_1,\vmui{1},\vmui{2}}\braces{\sum_{i=1}^{n_1}\normt{\vxii-\vmui{1}}+\sum_{i=n_1+1}^{n}\normt{\vxii-\vmui{2}}}~.
\end{align*}
The minimization on $\vmui{1},\vmui{2}$ is immediate and is given by the means of the segments, so we get
\begin{align}\label{eq:2seg_l0constraint}
    \argmin{n_1}\braces{\sum_{i=1}^{n_1}\normt{\vxii-\bar{x}_1}+\sum_{i=n_1+1}^{n}\normt{\vxii-\bar{x}_2}}~,
\end{align}
where we defined
\begin{align*}
\bar{x}_1&=\frac{1}{n_1}\sum_{i=1}^{n_1}\vxii\\
\bar{x}_2&=\frac{1}{n_2}\sum_{i=n_1+1}^{n}\vxii~
\end{align*}
and $n_2\triangleq n-n_1$.
Recall that the solution to the relaxed optimization problem is given by \lemref{lem:analyticalSolutionK=2}, as described in \secref{sec:outlierRobustObj}:
\begin{align}\label{eq:2seg_weighted_l1}
    \argmax{n_1,n_2=n-n_1}\braces{\frac{n_1 n_2}{nw_{n_1}}\vnorm{\bar{x}_2-\bar{x}_1}_2},
\end{align}
where $w_i$ are the weights. We now argue that \eqref{eq:2seg_l0constraint} and \eqref{eq:2seg_weighted_l1} have the same solution, for the specific choice of $w_i=\sqrt{i\paren{n-1}/n}$. First note that \eqref{eq:2seg_l0constraint} can be rewritten as $\argmax{n_1,n_2=n-n_1}\braces{n_1\normt{\bar{x}_1}+n_2\normt{\bar{x}_2}}$. Next, we show that this objective and the square of the (non-negative) objective \eqref{eq:2seg_weighted_l1} differ by a constant $C$, which depends on the data $\vxii$ but not on $n_1$:
\begin{align*}
    n_1\normt{\bar{x}_1}+n_2\normt{\bar{x}_2}&=\frac{n_1 n_2}{n}\normt{\bar{x}_2-\bar{x}_1}+C\\
    \frac{n_1^2}{n}\normt{\bar{x}_1}+\frac{n_2^2}{n}\normt{\bar{x}_2}&=C-\frac{2n_1 n_2}{n}\bar{x}_1^T\bar{x}_2\\
    \frac{1}{n}\normt{\sum_{i=1}^n\vxii}&=C,
\end{align*}
which proves that indeed  \eqref{eq:2seg_l0constraint} and \eqref{eq:2seg_weighted_l1} attain their optimal value at the same $n_1^*$
\end{proof}\QED
\newpage
\subsection{Proof of \eqref{eq:1D_errBound_m}}\label{sec:proof:1D_errBound_m}
\begin{proof}
Define the random variable
$Y_m$ which is the difference between the empirical means of two
subsequences created by splitting after $n_1+m$ samples:
\begin{align*}
    Y_m &\triangleq\frac{1}{n_2-m}\negspaces\sum_{i=n_1+m+1}^{n}\negspace\vxii-\frac{1}{n_1+m}\sum_{i=1}^{n_1+m}\vxii.
\end{align*}
This allows us to rewrite
\eqref{eq:lambdaStarWeighted} as an optimization over $m$:
\begin{align*}%\label{eq:optimal_m_star}
    %\!\!\!\!g\paren{m} &\!\triangleq\! \frac{(n_1\!+\!m)(n_2\!-\!m)}{n}\abs{Y_m},\\ %\quad,\quad
%    m^* \!=\!\argmax{m}{g\paren{m}},% \!\!\!%,%\nonumber
    g\paren{m} &\triangleq \frac{(n_1+m)(n_2-m)}{n}\abs{Y_m},\\
    m^* &=\argmax{m}\braces{g\paren{m}},
\end{align*}
where $m\in\brackets{-n_1+1,n_2-1}$.
Without loss of generality, we treat the case where
$m\geq0$. Note that $m^*\neq0$ if $g(0)<g(m)$ for some $m>0$.
%(it is not necessarily the optimal value $m^*$).
The probability of this event is:
\begin{align}\label{eq:correctBoundaryErrProb}
    &\prob{g(0) < g(m)}=\\
    &\prob{\frac{n_1n_2}{n}\abs{Y_0}<\frac{(n_1+m)(n_2-m)}{n}\abs{Y_m}}.\nonumber
\end{align}
Defining the following random variables:
\begin{equation}\label{eq:Wdef}
    W_m^{\pm}\triangleq\frac{n_1n_2}{n}Y_0\pm\frac{(n_1+m)(n_2-m)}{n}Y_m,
\end{equation}
we can rewrite
\begin{align}\label{eq:errBound_asW}
    \prob{g(0)<g(m)} & \leq\prob{W_m^{+}<0}+\prob{W_m^{-}<0},%\\
    %& \leq 2\prob{W_m^{-}<0} \leq 2\exp\paren{-Cm},\nonumber
\end{align}
where we used the fact that
%Our starting point is the bound given in \eqref{eq:correctBoundaryErrProb}, and the fact that
for two random variables $A$ and $B$, it holds true that
\begin{align*}%\label{eq:abProbBound}
    \prob{\abs{A}<\abs{B}}\leq\prob{A<B}+\prob{A<-B}.
\end{align*}
Using the definition of $Y_m$, we rewrite $W_m^{\pm}$
%, as defined by \eqref{eq:Wdef},
as the (weighted) average of a sequence of the $n$ random variables composing the data,
\begin{align}
    W_m^+=\frac{1}{n}\Bigg(&\paren{m-2n_2}\sum_{i=1}^{n_1}\vxii+\paren{n_1-n_2+m}\negspaces\sum_{i=n_1+1}^{n_1+m}\vxii\nonumber\\
    &+\paren{2n_1+m}\negspace\sum_{i=n_1+m+1}^{n}\negspace\vxii\Bigg),\nonumber
\end{align}
and
\[
   %W_m^-=-\frac{n_1m}{n}X_{1,n_1}+\frac{m(n-m)}{n}X_{n_1+1,n_1+m}+\frac{m(m-n_2)}{n}X_{n_1+m+1,n}
    W_m^-=\frac{1}{n}\paren{-m\sum_{i=1}^{n_1}\vxii+\paren{n-m}\negspaces\sum_{i=n_1+1}^{n_1+m}\negspaces\vxii-m\negspace
    \sum_{i=n_1+m+1}^{n}\negspace\vxii}.
\]
The means of $W_m^{\pm}$ are given by
\begin{align*}
    \expec{W_m^-} &= \frac{n_1m}{n}\Delta\mu,\\
    \expec{W_m^+} &= \frac{n_1\paren{2n_2-m}}{n}\Delta\mu,
\end{align*}
where we define $\Delta\mu\triangleq\mu_2-\mu_1$. From \eqref{eq:correctBoundaryErrProb}, \eqref{eq:Wdef}, and \eqref{eq:errBound_asW} it follows that
\begin{align*}
    &\prob{\frac{n_1n_2}{n}Y_0<\frac{(n_1+m)(n_2-m)}{n}\abs{Y_m}}\\
    \leq&\prob{W_m^-<0}+\prob{W_m^+<0}.\nonumber
\end{align*}
Using Hoeffding's inequality to bound the probabilities that $W_m^{\pm}$ are negative we get:
\begin{align*}%\label{eq:probWneg_hoeffding}
    \prob{W_m^-<0}\leq\exp\paren{-\frac{2n_1^2m\Delta\mu^2}{M^2n\paren{n-m}}}\triangleq B^-,
\end{align*}
and similarly:
\begin{align*}%\label{eq:probWpos_hoeffding}
    &\prob{W_m^+<0}\\
    &\leq\exp\paren{-\frac{2n_1^2\paren{2n_2-m}^2\Delta\mu^2}{M^2n\paren{m\paren{n-m-4n_1}+4n_1\paren{n-n_1}}}}\\
    &\triangleq B^+.
\end{align*}
%By comparing the arguments of the exponents in \eqref{eq:probWneg_hoeffding} and \eqref{eq:probWpos_hoeffding},
It can be shown that $B^+<B^-$, provided that $m$ is in its feasible range, i.e $0<m<n_2$. We conclude that
\begin{align*}
    &\mathbb{P}\paren{\frac{n_1n_2}{n}\abs{Y_0}<\frac{(n_1+m)(n_2-m)}{n}\abs{Y_m}}\\
    &\leq B^-+B^+\leq 2B^-,
\end{align*}
so we have that
\begin{align*}
    &\mathbb{P}\paren{\frac{n_1n_2}{n}\abs{Y_0}<\frac{(n_1+m)(n_2-m)}{n}\abs{Y_m}}\\
    &\leq2\exp\paren{-\frac{2n_1^2m\Delta\mu^2}{M^2n\paren{n-m}}}\leq2\exp\paren{-\frac{2n_1^2\Delta\mu^2}{M^2n^2}m},
\end{align*}
which proves \eqref{eq:1D_errBound_m} for
$C=\frac{2\Delta\mu^2n_1^2}{M^2n^2}$.
\end{proof} \QED
\newpage
\subsection{Proof of \thmref{thm:1D_errBound}}\label{sec:proof:1D_errBound}
\begin{proof}
Using \eqref{eq:1D_errBound_m} and the union bound, we get
\begin{align*}
    \prob{m^*\geq m_0} &\leq \prob{\exists m\geq m_0:\:g(0)<g(m)} \\
    &\leq\sum_{m=m_0}^{n_2}\prob{g(0)<g(m)} \\
    &\leq n_2\prob{g\paren{0}<g\paren{m_0}} \\
    &\leq 2n_2\exp\paren{-Cm_0} \stackrel{(\star)}{\leq} \delta~,
\end{align*}
where ($\star$) follows from %The last inequality follows %from
the definition of $m_0$.
\end{proof}\QED

\subsection{Proof of \lemref{lem:replaceSamplesWithMean}}\label{sec:proof_lem:replaceSamplesWithMean}
\begin{proof} Consider a segment of $n_0$ data samples, and denote by $\vmui{0}$ the centroid of this segment. The mean
of the segment is given by $\bvxi{0}=\frac{1}{n_0}\sum\limits_{\vxii\in\vmui{0}}\vxii$. The contribution of this segment to
the first term in the objective \eqref{eq:segmentationObjective_mu} is given by
\begin{equation}\label{eq:oneSegContribution}
\frac{1}{2}\sum\limits_{\vxii\in\vmui{0}}\negspaces\normt{\vxii-\vmui{0}}=
\frac{1}{2}\sum\limits_{\vxii\in\vmui{0}}\negspaces\left(\normt{\vxii}-2\vxii^T\vmui{0}+\normt{\vmu_0}\right).
\end{equation}
If we now substitute $\bvxi{0}$ for each sample $\vxii$ in this segment, the contribution to the objective becomes
\begin{equation}\label{eq:oneSegMeanContribution}
    \frac{n_0}{2}\normt{\bvxi{0}-\vmui{0}}=\frac{n_0}{2}\left(\normt{\bvxi{0}}-2\bvxi{0}^T\vmui{0}+\normt{\vmui{0}}\right).
\end{equation}
It is straightforward to show that as a function of $\vmui{0}$, \eqref{eq:oneSegContribution} and
\eqref{eq:oneSegMeanContribution} differ by a constant which depends only on the data samples $\vxii\in\vmui{0}$,
not on $\vmui{0}$. Since this argument holds for all segments, we conclude that replacing each data sample with the mean
of the segment to which it belongs, results in the same objective, up to a constant. Therefore the optimal solution $\vmu^*$
does not change.
\end{proof}\QED
